%% kapa: exam=1 round=35 year=2024 period=1 form=A subject=economics subject_ko=경제학원론 seq=10 number=50
%% source: Q-Net 감정평가사 기출문제 (cst00309) artlSeq=5213229 file=raw/1008/5213229/1교시 문제지 원본.pdf
%% conversion: kordoc 4.24.51 -> tools/gen1.py (auto); review=proofread-sonnet; flags=none
\begin{question}{50}
독점기업의 독점력과 가격규제 정책에 관한 설명으로 옳지 않은 것은?
\begin{choices}
  \item 러너의 독점력지수(Lerner index of monopoly power)는 수요곡선 상의 이윤극대화점에서 측정한 수요의 가격탄력성의 역수와 같은 값이다.
  \item 한계비용가격설정은 자연독점 기업에게 손실을 초래한다.
  \item 평균비용가격설정은 기업이 손실을 보지 않으면서 가능한 많은 상품을 낮은 가격에 공급하도록 유도할 수 있다.
  \item 이중가격설정(two-tier pricing)은 한계비용가격설정의 장점을 살리면서도 독점기업의 손실을 줄일 수 있도록 하는 정책이다.
  \item 이중가격설정은, 낮은 가격은 한계비용과 한계수입이 일치하는 가격으로, 높은 가격은 한계비용곡선과 수요곡선이 교차하는 지점의 가격으로 판매하도록 하는 정책이다.
\end{choices}
\end{question}
\answer{5}
